\chapter{CONCLUSION AND FUTURE SCOPE}

\section{CONCLUSION}
The ScrapKart AI project was conceptualized and developed to resolve the fundamental inefficiencies and lack of transparency plaguing the traditional scrap collection and recycling industry. Through the rigorous integration of Artificial Intelligence and modern web technologies, this project successfully transformed a highly informal sector into a structured, data-driven digital platform. 

The implementation of the YOLOv8-based computer vision model proved highly effective in automating the identification and initial valuation of recyclable materials from user-uploaded images. By coupling this visual analysis with a dynamic pricing engine linked to live market data, ScrapKart AI established a transparent and fair economic mechanism, directly addressing consumer hesitations regarding fair compensation. Furthermore, the integration of advanced route optimization algorithms provided tangible operational benefits to scrap collectors, significantly reducing their travel distance, saving fuel, and increasing their daily collection efficiency. 

From a broader perspective, ScrapKart AI serves as a powerful tool for environmental conservation. The inclusion of a sustainability dashboard successfully incentivized users by quantifying their contributions to carbon footprint reduction. In conclusion, ScrapKart AI not only streamlines logistics and ensures economic fairness but also actively promotes a sustainable, circular economy, proving that AI-driven solutions can successfully tackle complex socio-environmental challenges at the grassroots level.

\sectionsection{FUTURE SCOPE}
While ScrapKart AI represents a significant advancement over traditional methods, there are several avenues for future enhancement and expansion:
\begin{itemize}
    \item \textbf{Integration of Hardware IoT Scales:} To overcome the limitation of estimating exact weight from 2D images, future iterations could integrate Bluetooth-enabled smart weighing scales used by collectors. These scales would automatically transmit the precise weight to the application upon pickup, recalculating the final price with 100\% accuracy.
    \item \textbf{Expansion to Industrial Waste Management:} The current platform is optimized for household B2C (Business-to-Consumer) transactions. The architecture can be scaled to handle B2B (Business-to-Business) operations, facilitating the bulk sale of industrial scrap metal and electronic waste from factories.
    \item \textbf{Advanced AI Models for Material Grading:} Beyond simply identifying a material as "Plastic," future AI models could be trained to identify specific polymer grades (e.g., PET vs. HDPE), which have vastly different market values and recycling requirements.
    \item \textbf{Government and Municipal Integration:} Collaborating with local municipalities could allow the platform to serve as the official digital infrastructure for a city's waste management grid, enabling authorities to track city-wide recycling metrics and identify areas requiring improved waste disposal awareness.
\end{itemize}
